\documentclass[12pt,a4paper]{article}
\usepackage[utf8]{inputenc}
\usepackage[T1]{fontenc}
\usepackage{lmodern}
\usepackage{amsmath,amssymb}
\usepackage{siunitx}
\usepackage{geometry}
\usepackage{hyperref}
\usepackage{enumitem}
\usepackage{tikz}
\geometry{margin=2.2cm}

\title{Lecture 3 -- Frequency Response and Bandwidth}
\author{Readable Assignment Set}
\date{}

\begin{document}
\maketitle
\tableofcontents

\section{Assignments}
\subsection{C1}
Op-amp with unity-gain bandwidth \SI{15}{MHz}. Used as non-inverting amplifier with DC gain 10 and then 100. Find 3-dB bandwidth in both cases.

\subsection{C2}
Dominant-pole op-amp with $A_{0,OL}=\num{50000}$ and $f_{BOL}=\SI{100}{Hz}$. Sketch magnitude Bode plot of $A_{OL}$ to scale.

\subsection{C3}
Op-amp with $A_{0,OL}=\num{200000}$ and 3-dB bandwidth \SI{5}{Hz}. Find magnitude and phase of open-loop gain at \SI{100}{Hz}, \SI{1000}{Hz}, and \SI{1}{MHz}.

\subsection{C4}
Design gain 100 using either (1) one non-inverting stage of gain 100 or (2) two cascaded stages of gain 10. Op-amp GBP is $10^6$. Write gain vs frequency and find 3-dB bandwidth for each option.

\subsection{C5}
Dominant-pole op-amp with DC gain 200,000 and 3-dB bandwidth \SI{5}{Hz}. Sketch to scale the Bode magnitude of open-loop gain. Then sketch closed-loop magnitude for non-inverting gains 100 and 10.

\subsection{C6}
Inverting amplifier with $R_2=\SI{100}{k\Omega}$, $R_1=\SI{10}{k\Omega}$, input 
\[
v_{in}=0.1\sin(2\pi\cdot 5\cdot10^3 t)\,\si{V},
\]
with open-loop model
\[
A_{OL}(f)=\frac{A_0}{1+jf/f_c},\quad A_0=100\,\text{dB},\ f_c=\SI{10}{Hz}.
\]
1) Find output amplitude assuming ideal op-amp.
2) Find percentage error from finite open-loop behavior.

\subsection{C7}
Non-inverting amplifier with $R_2=\SI{99}{k\Omega}$, $R_1=\SI{1}{k\Omega}$ and same open-loop model as C6.
\begin{enumerate}
\item Sketch Bode magnitude of $A_{OL}$, $1/\beta$, and $A_V$ from \SI{1}{Hz} to \SI{10}{MHz}.
\item Find 3-dB bandwidth $f_1'$.
\item Derive symbolic expression for $f_1'$ as function of $\beta$ and open-loop parameters.
\end{enumerate}

\section{Hints}
\begin{itemize}
\item For dominant-pole op-amp: $A_{OL}(j f)=\dfrac{A_0}{1+j f/f_c}$.
\item GBP relation: $f_t\approx A_{CL}(0)\,f_{BCL}$.
\item Closed-loop transfer for feedback amplifiers: 
$A_{CL}=\dfrac{A_{ideal}}{1+\dfrac{1}{\beta A_{OL}}}$.
\item For phase of a first-order pole: $\angle A_{OL}=-\arctan(f/f_c)$.
\end{itemize}

\section{Step-by-step solutions}
\subsection{C1}
\[
f_{3dB}\approx \frac{f_t}{A_{CL}(0)}.
\]
For gain 10: $f_{3dB}=15\,\text{MHz}/10=\SI{1.5}{MHz}$.
For gain 100: $f_{3dB}=15\,\text{MHz}/100=\SI{150}{kHz}$.

\subsection{C2}
\begin{enumerate}
\item $20\log_{10}(A_0)=20\log_{10}(50000)=\SI{93.98}{dB}$.
\item Break at $f_c=\SI{100}{Hz}$.
\item Slope above break: \SI{-20}{dB/dec}.
\item Unity crossing near $f_t=A_0f_c=\SI{5}{MHz}$.
\end{enumerate}

\subsection{C3}
Given $A_0=200000$ and $f_c=\SI{5}{Hz}$:
\[
|A_{OL}|=\frac{A_0}{\sqrt{1+(f/f_c)^2}},\quad \angle A_{OL}=-\arctan(f/f_c).
\]
\begin{itemize}
\item $f=\SI{100}{Hz}$: $f/f_c=20$, $|A|\approx 9990$ (\SI{79.99}{dB}), phase $\approx \SI{-87.1}{\degree}$.
\item $f=\SI{1000}{Hz}$: $f/f_c=200$, $|A|\approx1000$ (\SI{60.00}{dB}), phase $\approx \SI{-89.7}{\degree}$.
\item $f=\SI{1}{MHz}$: $f/f_c=200000$, $|A|\approx1$ (\SI{0}{dB}), phase $\approx \SI{-90}{\degree}$.
\end{itemize}

\subsection{C4}
\textbf{Alternative 1: one stage, gain 100}
\[
A_1(jf)=\frac{100}{1+jf/f_{B1}},\quad f_{B1}=\frac{10^6}{100}=\SI{10}{kHz}.
\]

\textbf{Alternative 2: two stages, each gain 10}
\[
A_2(jf)=\left(\frac{10}{1+jf/f_{B2}}\right)^2,\quad f_{B2}=\frac{10^6}{10}=\SI{100}{kHz}.
\]
Overall DC gain is 100. The composite 3-dB point satisfies
\[
\left(1+(f/f_{B2})^2\right)^2=2 \Rightarrow \frac{f}{f_{B2}}=\sqrt{\sqrt{2}-1}\approx0.6436
\Rightarrow f\approx0.6436f_{B2}=\SI{64.4}{kHz}.
\]
So cascaded design gives much larger bandwidth than single-stage \SI{10}{kHz}.

\subsection{C5}
Open-loop: same as C3, with DC \SI{106}{dB} and break \SI{5}{Hz}, slope \SI{-20}{dB/dec}.
Closed-loop gain 100: flat at \SI{40}{dB} until about \SI{10}{kHz}.
Closed-loop gain 10: flat at \SI{20}{dB} until about \SI{100}{kHz}.

\subsection{C6}
\begin{enumerate}
\item Ideal inverting gain magnitude $|A|=R_2/R_1=10$.
Input amplitude is \SI{0.1}{V} so output amplitude is \SI{1.0}{V}.
\item At \SI{5}{kHz}: 
\[
|A_{OL}|\approx \frac{10^{100/20}}{\sqrt{1+(5000/10)^2}}\approx\frac{100000}{500}\approx200.
\]
Feedback factor for inverting stage with grounded source: 
$\beta=\frac{R_1}{R_1+R_2}=\frac{1}{11}$.
Loop gain $\beta|A_{OL}|\approx18.18$.
Relative gain error estimate $\approx\frac{1}{1+\beta|A_{OL}|}\approx\SI{5.2}{\percent}$.
\end{enumerate}

\subsection{C7}
\begin{enumerate}
\item $\beta=\frac{R_1}{R_1+R_2}=\frac{1}{100}$, so $1/\beta=100$ (\SI{40}{dB}).
\item With $f_t=A_0f_c=100000\cdot10=\SI{1}{MHz}$ and non-inverting gain 100, 
\[
f_1'\approx\frac{f_t}{100}=\SI{10}{kHz}.
\]
\item Symbolic expression using single-pole model:
\[
f_1'\approx\beta A_0 f_c = \beta f_t,
\]
valid when $A_0\gg1$ and single dominant pole approximation is valid.
\end{enumerate}

\subsection{Illustrative Bode sketch}
\begin{center}
\begin{tikzpicture}[xscale=1.1,yscale=0.09]
\draw[->] (0,0) -- (10,0) node[right] {$\log f$};
\draw[->] (0,0) -- (0,115) node[above] {dB};
\draw[thick] (0,94) -- (3,94) -- (9, -26);
\draw[dashed] (3,0) -- (3,94) node[above] {$f_c$};
\draw[dashed] (7.7,0) -- (7.7,10) node[above] {$f_t$};
\end{tikzpicture}
\end{center}

\end{document}
